% 致谢
\chapter*{致\hspace{1em}谢}
\addcontentsline{toc}{chapter}{致谢}
\markboth{致谢}{致谢}

本论文的完成离不开学校支持以及老师们与朋友的帮助。

在完成论文的最后，我首要要感谢指导教师金蒙豪老师。自大一 SRTP 课题互选阶段起，金老师就引导我进入微磁学的研究领域，最开始指导我研究基础的磁性拓扑结构，也就是二维斯格明子，并一起发表了相关的专利。后来又和我们一起讨论研究霍普夫子的相关研究，使我的研究对象从二维斯格明子逐步深入至三维霍普夫子。在本次毕业设计期间，从课题选定、研究方案设计到最终论文的修改定稿，金老师都提供了耐心的学术指导，为我在磁性拓扑结构方向的研究奠定了基础。

此外，本课题的Mumax3仿真工作离不开杭州电子科技大学电子信息学院 519 实验室提供的计算资源与科研平台。在初期缺乏算力的情况下，实验室的支持保障了微磁学仿真计算的顺利推进。同时需要感谢同组同学在日常学术讨论、组内汇报以及代码调试环节中提供的协助，在最开始研究斯格明子的时候，一起从空白领域学习，慢慢一步一个脚印的学会一些基础的微磁学仿真内容和OOMMF、Mumax3等工具的使用。

在本科四年的学习过程中，各位授课教师的专业教导以及班级好友、寝室室友的互助陪伴同样必不可少。感谢同学们在日常学习与生活中的鼓励。最后，家人的理解与包容是我顺利完成本科阶段学业的重要后盾，在此向他们表达诚挚的谢意。谨以此文感谢所有曾给予我帮助的师长与朋友，并致谢学校四年的指导和支持。
